\documentclass[10pt]{article}
\usepackage[margin=20mm]{geometry}
\usepackage{amsmath,amssymb,mathtools,fontspec,unicode-math,xurl,hyperref}
\IfFontExistsTF{DejaVu Serif}{\setmainfont{DejaVu Serif}}{\setmainfont{Latin Modern Roman}}
\setmathfont{latinmodern-math.otf}
\setlength{\parindent}{0pt}\setlength{\parskip}{6pt}\sloppy
\title{Technical Review: Finite-Energy Reparametrization for Algebraic Hypersurfaces}
\author{AI-assisted technical review}
\date{}
\begin{document}
\maketitle
Abstract: This review checks the finite-energy reparametrization argument in the publisher version of Optimal Transport Between Algebraic Hypersurfaces. A continuous semialgebraic family of binary quadratic forms satisfies every stated hypothesis of Theorem 5.13, but every asserted power reparametrization has infinite quadratic Wasserstein energy. The obstruction is a parameter-speed singularity at a regular terminal hypersurface. This is a Major defect in a lemma used in the core metric construction. It does not refute the final existence theorem. Scientific adjudication and independent acceptance remain pending.\par
\section*{Identity, reviewed version, and scope}
The reviewed paper is Paolo Antonini, Fabio Cavalletti, and Antonio Lerario, Optimal Transport Between Algebraic Hypersurfaces, Geometric and Functional Analysis 35 (2025), 43–112, DOI 10.1007/s00039-025-00699-w (https://doi.org/10.1007/s00039-025-00699-w) [ACL]. The assigned identity is M25-0371. The source is the acquired 70-page publisher version of record, bound to the hashes in Appendix 5 (\#app:record). Throughout this report, source page references are printed journal pages; the corresponding one-based physical PDF page is the printed page minus 42.\par
The focused review covers the definitions of admissibility and energy, the quantifiers in Theorem 5.13, its proof, and its explicit use in Corollary 5.14 and Theorem 5.15. The one-dimensional measure construction and the paper’s Fubini–Study normalization are checked where needed for the counterexample. The continuity/injectivity discussion and the lower-semicontinuity argument received a supporting text inspection, not a complete independent reconstruction. The general Hamiltonian isotopy argument, the Kähler calculation, the higher-dimensional root-regularity results, and the condition-number application have not been fully verified. This is a claim-focused initial review, not a certificate of whole-paper correctness or independent acceptance.\par
\section*{Definitions and the precise claim}
Let \(P_{n,d}\) denote the projective space of nonzero degree-\(d\) homogeneous polynomials in \(n+1\) complex variables, and let \(\Delta_{n,d}\) be the discriminant of polynomials whose projective zero equation is singular. In dimension \(n=1\), the measure \(\mu([p])\) is the probability measure assigning to each projective root its multiplicity divided by \(d\); see pp. 45–47 and Theorem 3.8, p. 69.\par
Definition 5.1, p. 77, requires \[\mu\circ\gamma\in AC^2([0,1];\mathcal P_2(\mathbb{CP}^{n})),
 \qquad
 \mathcal E(\gamma)=\int_0^1|\dot\mu_t|^2\,dt.\] The speed is measured using the ambient \(W_2\) distance. The formal definition of \(AC^2\) in Definition 2.2, p. 57, supplies an \(L^2\) upper gradient for distances between every two times. Thus an almost-everywhere metric speed with divergent square integral rules out admissibility. References below to infinite energy use the extended-energy convention; they do not assert that a nonadmissible curve belongs to the domain of Definition 5.1.\par
Problem 1 (A terminal speed singularity survives every asserted power reparametrization).  Theorem 5.13, p. 86 (PDF p. 44), asserts that whenever a continuous semialgebraic curve \(\gamma:[0,1]\to P_{n,d}\) satisfies \[\gamma(0)\in\Delta_{n,d},\qquad
 \gamma(t)\notin\Delta_{n,d}\quad(0<t\leq1),\] there is a positive integer \(m\) for which \(\alpha_m(s)=\gamma(s^m)\) has finite energy. This statement is false with its stated regularity hypothesis.\par
\subsection*{Every hypothesis is satisfied}
Equation reference: eq:example. Take \(n=1\), \(d=2\), and define \[
 r(t)=t\bigl(1+\sqrt{1-t}\bigr),\qquad
 p_t(z_0,z_1)=z_0^2-r(t)^2z_1^2,\qquad
 \gamma(t)=[p_t],\quad 0\leq t\leq1.\] The square root is its nonnegative real branch. The coefficient of \(z_0^2\) is always one, so the projective representative never vanishes. The curve is continuous and semialgebraic: its coefficients are obtained by polynomial operations from the algebraic branch \(u\geq0\), \(u^2=1-t\). In fact, it is smooth on \([0,1)\).\par
Equation reference: eq:measure. At \(t=0\), \(p_0=z_0^2\) has a double root at \([0:1]\), so \(\gamma(0)\in\Delta_{1,2}\). For every \(t>0\), \(r(t)>0\), and the two roots \[z_+(t)=[r(t):1],\qquad z_-(t)=[-r(t):1]\] are distinct and simple. The gradient of \(p_t\) at either affine root has nonzero first component \(\pm2r(t)\). This checks avoidance of the discriminant for all \(t\in(0,1]\), including \(t=1\), where \(r(1)=1\). Consequently \[
 \mu_t=\tfrac12\delta_{z_+(t)}+\tfrac12\delta_{z_-(t)}
 \quad(t>0),\qquad \mu_0=\delta_{[0:1]}.\] No additional coefficient smoothness at the terminal endpoint is assumed in the printed theorem.\par
\subsection*{The actual metric speed}
Appendix A.1, p. 106, normalizes the Fubini–Study metric by the Riemannian submersion from the unit sphere. For the real coordinate \(r\), the normalized lift \[b(r)=\frac{(r,1)}{\sqrt{1+r^2}}\] has horizontal derivative of squared norm \((1+r^2)^{-2}\). Hence the projective speed of either root is \(|r'(t)|/(1+r(t)^2)\).\par
Equation reference: eq:speed. Near any fixed \(t>0\), the two roots are separated by a positive distance. The uncrossed matching between \(\mu_t\) and \(\mu_{t+h}\) has cost tending to zero, whereas the crossed matching has a positive limiting cost. The transport problem between these two equal-weight two-point measures minimizes a linear function over the \(2\times2\) doubly stochastic matrices, whose extreme points are precisely these two matchings. Thus the uncrossed matching is optimal for sufficiently small \(h\). Both matched roots have the same speed, so averaging their squared speeds gives \[
 |\dot\mu_t|^2=g(t):=\frac{r'(t)^2}{(1+r(t)^2)^2},\qquad 0<t<1.\] This derivation uses the measure itself and is unaffected by rescaling a homogeneous representative.\par
Equation reference: eq:limit-original. Put \(u=\sqrt{1-t}\). Direct differentiation gives \[r(t)=1+u-u^2-u^3,\qquad
 r'(t)=-\frac{1-2u-3u^2}{2u}.\] It follows exactly that \[
 \lim_{t\uparrow1}(1-t)g(t)
 =\lim_{u\downarrow0}
 \frac{(1-2u-3u^2)^2}
 {4\{1+(1+u-u^2-u^3)^2\}^{2}}
 =\frac1{16}.\]\par
\subsection*{Failure for every positive integer \(m\)}
Equation reference: eq:limit-power. For \(\nu_s=\mu_{s^m}\), the local matching argument and the chain rule yield \[|\dot\nu_s|^2=m^2s^{2m-2}g(s^m).\] Since \[\frac{1-s^m}{1-s}=\sum_{j=0}^{m-1}s^j\longrightarrow m,\] equation [eq:limit-original] (\#eq:limit-original) gives, for each integer \(m\geq1\), \[
 \lim_{s\uparrow1}(1-s)|\dot\nu_s|^2=\frac{m}{16}>0.\] In particular, on a sufficiently small left neighborhood of \(1\), \[|\dot\nu_s|^2\geq\frac{m}{32(1-s)}.\] Therefore \(\int_0^1|\dot\nu_s|^2\,ds=\infty\), and \(\mu\circ\alpha_m\notin AC^2\), for every allowed \(m\). The asserted finite-energy conclusion is contradicted by a universal analytic calculation, not by testing finitely many integers.\par
\subsection*{The failing proof step and alternative readings}
In the proof on p. 87 (PDF p. 45), immediately after splitting the energy integral at \(r_0\), regularity of the hypersurface equation is used to conclude that the integrand is bounded on \([r_0,1]\). Regularity controls the geometric denominator but does not bound the derivative of a merely continuous semialgebraic coefficient path. In [eq:example] (\#eq:example), the roots stay separated near \(1\), while \(r'(t)\) has a square-root singularity and the square speed is nonintegrable.\par
The assertion and this proof step were checked in the actual publisher-page images; the issue is not inferred from an OCR anomaly. An additional \(C^1\) assumption up to the terminal endpoint would exclude this example. The source does impose \(C^1\) hypotheses in other results, including Theorem 5.3, but it does not impose that hypothesis in Theorem 5.13. A local assertion on a sufficiently short interval adjacent to \(0\) is also not refuted. Neither alternative is the global statement printed in the reviewed version.\par
\section*{Core dependency and repair boundaries}
Corollary 5.14, p. 87, invokes Theorem 5.13 to connect a point on the discriminant to a regular point by a finite-energy curve. The proof of Theorem 5.15, pp. 87–88, then uses Corollary 5.14 in establishing the intrinsic metric and its geodesic construction. The false quantified lemma and its invalid boundedness step therefore occur in an explicit proof chain supporting a core contribution. The supported impact is Major under the assigned core-derivation standard.\par
This conclusion has a precise boundary. The counterexample does not show that the two endpoints admit no finite-energy connection. It does not refute the continuity/injectivity of the measure map, the ultimate completeness or geodesicity claim, or the separate Kähler and condition-number results. The existence argument may be repairable by choosing better parameterized paths. No difficulty-of-repair claim is needed for the impact assessment.\par
Equation reference: eq:repair, eq:repair-bound. For this particular example there is an explicit verified repair using a different reparametrization: \[
 \phi(s)=1-(1-s)^2,\qquad
 r(\phi(s))=(2s-s^2)(2-s)=4s-4s^2+s^3 .\] The function \(\phi\) is a continuous, strictly increasing semialgebraic bijection of \([0,1]\) with its endpoints fixed. The repaired roots are smooth. Their \(W_2\) energy satisfies \[
 \mathcal E(\gamma\circ\phi)
 =\int_0^1\frac{(4-8s+3s^2)^2}
 {\{1+(4s-4s^2+s^3)^2\}^{2}}\,ds
 \leq\int_0^1(4-8s+3s^2)^2\,ds
 =\frac{47}{15}<\infty .\] The endpoint collision at \(s=0\) is harmless: pairing the two smooth roots supplies a bounded metric upper gradient on the whole interval. The quantity \(47/15\) is an upper bound, not the claimed exact energy. This verifies a repair of the example; it does not prove a revised global theorem for every continuous semialgebraic path or certify a complete replacement proof of Corollary 5.14.\par
\section*{Version checks and conclusion}
The inspected publisher PDF is the only full article version used for the finding. The public publisher HTML also displays the same unrestricted Theorem 5.13. A bounded exact-title correction/erratum search and a search of that article page identified no correction. The arXiv landing page lists v3, revised 16 July 2023, as its latest version [arxivmeta]; its full PDF was not substituted for the publisher version or compared line by line. These checks do not establish novelty of the finding or absence of an unpublished correction.\par
The acquisition manifest records bounded searches for supplements and research code without identifying a target-specific release. No author code was executed. No additional experiment is required to establish the analytic counterexample.\par
Assessment | Repair verification | Author adoption | Remaining work\par
Major, Problem 1 | Example reparametrization and finite-energy bound verified; general revision not certified | Unknown; not requested | Independent acceptance; corrected global statement and downstream proof review\par
Scientific decision: pending. Human reviews: none recorded. Initial delivery does not constitute independent acceptance.\par
\section*{Supporting checks}
The source PDF SHA256 is \[\text{\texttt{\scriptsize 9e183c6543a7edf531e15709eece1ab62a13e4fb60b69b8aead3e31480840fa6}}.\] The source-manifest SHA256 is \[\text{\texttt{\scriptsize 8c534e5a0f1ecd13a428d95f3841ced417ec6812583d667eb2d69097b30cb685}}.\] The accompanying `verification/0371` directory retains the source-page bindings, actual viewing receipts, the independent `proof\_check.py`, its execution result, and build/dependency records. The exact-arithmetic script verifies the endpoint coefficient \(1/16\), the polynomial identity in [eq:repair] (\#eq:repair), and the bound \(47/15\). Its finite checks of selected \(m\) values are algebra diagnostics; the argument for every positive integer \(m\) is [eq:limit-power] (\#eq:limit-power).\par
The scientific analysis and image inspections are performed in the assigned current Chat. The controller-supplied workspace/model UI observation is attributed to that controller; no independent backend-equivalence assertion or human attestation is made. Complete source and execution limitations are retained in the coverage and provenance records.\par
9 P. Antonini, F. Cavalletti, and A. Lerario, Optimal Transport Between Algebraic Hypersurfaces, Geom. Funct. Anal. 35 (2025), 43–112. Publisher version of record. https://doi.org/10.1007/s00039-025-00699-w. Public article HTML consulted for the bounded correction and statement check.\par
P. Antonini, F. Cavalletti, and A. Lerario, Optimal transport between algebraic hypersurfaces, arXiv:2212.10274, v3 (16 July 2023). Landing-page metadata and revision history only. https://arxiv.org/abs/2212.10274.\par
\end{document}
